\section{\centering Fixed Point Combinators and Recursion}
\label{sec:recursion-fixed-points}

The Church Rosser Property labels the untyped \lcalc \ as a confluent system, therefore we know that whenever a term has a normal form it must be unique. We still have not given an answer to whether all \lterms \ have a normal form.

The most straight forward way to test this, is to check whether there exists some term that reduces to itself, effectively creating a loop in the $\beta$-reduction chain so to speak. In more formal terms this means there exists some $M$ such that \( M \curly^n M \) with $ n > 0 $. It turns out \( \Omega \equiv (\la x . \, x \, x) (\la x . \, x \, x) \) fullfills this requirement. Let us inspect what goes on while $\beta$-reducing this term:
\[
  \Omega \equiv (\la x . \, x \, x) (\la x . \, x \, x) \curly (\la x . \, x \, x) (\la x . \, x \, x) \curly^* (\la x . \, x \, x) (\la x . \, x \, x) \equiv \Omega
\]
This so called $\Omega$-combinator, proves that not all terms have a normal form, and shines a beam of light onto how recursion can be achieved in the \lcalc.

There is one caveat to this, that classically, recursion is achieved through naming, take for instance this definition for the $\mathrm{gdc}$:
\begin{align*}
  \underline{\mathrm{gcd}}(a,b) &= 
                      \begin{cases}
                        a & b = 0 \\
                        \underline{\mathrm{gcd}}(b,\, a \bmod b) & b \neq 0
                      \end{cases}
\end{align*}

Here, when $b \neq 0$ the name gcd is used to refer to the same definition again. In the \lcalc \ such self-referencing definitions are not permited since the language does not allow names.

\subsubsection{\centering Fixed Points}
To bypass this setback the notion of fixed point is of particular interest.
\begin{definition}
  Let \( f : X \to X \) be a function. An element \( x \in X \) is called a fixed point of \( f \) if \( f(x) = x. \)
\end{definition}
\begin{example} Some simple functions that illustrate what fixed points are:
  \[
    \begin{tabular}{|c|c|}
      \hline
      \textit{Function} & \textit{Fixed Point(s)} \\
      \hline
      $f(x) = c$ & $x = c$ \\
      \hline
      $g(x) = x$ & $\{ x \mid x \in \mathbb{R} \}$ \\
      \hline
      $h(x) = x^2$ & $ x = 0,\, x = 1 $ \\
      \hline
    \end{tabular}
  \]
\end{example}
This harmless idea satisfies an interesting property, namely that $ x = f(x) = f(f(x)) = f^n(x) $ effectively creating a chain where every element $ f^{n+1}(x) = f(f^{n}(x))$, attaining self-reference without naming and showcasing the ability of fixed points to induce recursion. Our quest now, is to find a term that simulates this behavior. To do this, let us translate the definition for fixed point into the world of \lcalc.
\begin{definition}
  Given a lambda term \( L \in \La \), a term \( M \in \La \) is called a fixed point of \( L \) if \( M =_\beta L M \).
\end{definition}
No surprises here, the definition is close to a one-to-one mapping of its functional counterpart seen previously. Interestingly enough, it turns out that every lambda term has a fixed point:
\begin{theorem}
  For every \( L \in \La \), there exists \( M \in \La \) such that $ M $ is a fixed point for $ M $ i.e. $ M =_\beta L M $. Referred to as the Fixed Point Theorem in the literature.
\end{theorem}
\begin{proof} Given $L \in \La$, define $M$ as below:
  \[
    M := (\la x . \, L \, (x \, x))(\la x . \, L \, (x \, x))
  \]
  $M$ is a reducible expression:
  \begin{align*}
    M 
    &\equiv (\la x . \, L \, (x \, x))(\la x . \, L \, (x \, x)) \\
    &\curly L \, ((\la x . \, L \, (x \, x))(\la x . \, L\, (x \, x))) \\
    &\equiv L \, M.
  \end{align*}
  Hence, $ M =_\beta L M $, as required.
\end{proof}
\begin{remark}
  See how $ L M $ is the application of $ M $ to $ L $ and that using \textsc{AppR} we can keep $\beta$-reducing $ M $ the same way we did for $ x = f(x) = f(f(x))\dots $ only this time using lambda terms:
  \[
    M \curly L \, M \curly L \, L \, M \curly^* L^n M.
  \]
\end{remark}
\subsubsection{\centering Recursive Combinators}
Now that nameless self-reference has been proven to exist within the the \lcalc, we take a step forward and introduce a few of the most important fixed point combinators and explore their different natures.
\begin{definition} $Y$ Combinator. This is the term that given some other term returns its fixed point.
\[
  Y \triangleq \la f.\, (\la x.\, f\, (x\, x))\, (\la x.\, f\, (x\, x))
\]
\end{definition}
\begin{remark}
  Even though the $Y$ combinator is of theoretical interest since it allows us to find recursion within the \lcalc, it is of little practical under strict evaluation since it always runs into non-termination e.g.
  \[
    \begin{aligned}
      Y \, F &\equiv (\la f. \, (\la x . \, f \, (x \, x))(\la x . \, f \, (x \, x))) F \\
      &\curly (\la x . \, F \, (x \, x))(\la x . \, F \, (x \, x)) \\
             &\curly F((\la x . \, F \, (x \, x))(\la x . \, F \, (x \, x))) \\
             &\curly F(F((\la x . \, F \, (x \, x))(\la x . \, F \, (x \, x)))) \\
             &\curly \cdots
    \end{aligned}
  \]
  Under strict evaluation, the argument $F$ is not evaluated until $Y$ is done unraveling i.e. never.
\end{remark}
To work around this issue, several other combinators exist, the $Z$ and the Turing combinators are two of the better known ones:
\begin{definition} \( Z \triangleq \la f.\, (\la x.\, f\, (\la v.\, x\, x\, v))\, (\la x.\, f\, (\la v.\, x\, x\, v)) \)
\end{definition}
\begin{example} Unfolding the $Z$ Combinator
  \[
    \begin{aligned}
      Z \, F &\equiv (\la f.\, (\la x.\, f\, (\la v.\, x\, x\, v))\, (\la x.\, f\, (\la v.\, x\, x\, v)))\, F \\
             &\curly (\la x.\, F\, (\la v.\, x\, x\, v))\, (\la x.\, F\, (\la v.\, x\, x\, v)) \\
             &\curly F\, (\la v.\, (\la x.\, F\, (\la v.\, x\, x\, v))\, (\la x.\, F\, (\la v.\, x\, x\, v))\, v)\, \\
             &\equiv F\, (\la v.\, Z \, F \, v)\, \\
    \end{aligned}
  \]
\end{example}
\begin{remark}
  The extra $\lambda v.\,\cdot$ in the final form of $Z\,F$ is essential to delay evaluation in a strict environment. Without it, the recursive call would be applied immediately, and cause non-termination in a call-by-value setting.
\end{remark}
\begin{definition} Turing combinator \( \Theta \triangleq (\lambda x.\, \lambda y.\, y\, (x\, x\, y))\, (\lambda x.\, \lambda y.\, y\, (x\, x\, y)) \)
\end{definition}
We now use the Turing combinator to aid us implement an untyped version of the factorial function.
\subsubsection{\centering \lCalc \ Implementation of Factorial}
Classically, factorial is defined as $n! = \prod_{k=1}^{n} k$, that has the following recursive interpretation $n! = n \cdot (n-1)!$ with $0! = 1$. This section develops a worked out example of how to implement this in the \lcalc.
\begin{example} Implementing factorial using the acquired toolkit:
  \[
    \underline{\texttt{fact}}
    = \lambda n.\; \texttt{if\ then\ else}\;(\texttt{iszero}\; n)\;(\overline{\texttt{1}})\;(\texttt{mult}\; n\;(\underline{\texttt{fact}}\;(\texttt{pred}\; n))),
  \]
\end{example}
\noindent As already seen, the problem with a recursive definition such as this one is that the name $\texttt{fact}$, appears both on the left and the right-hand side. This means finding $\texttt{fact}$ requires solving an equation, which we will do thanks to fixed points. We start by abstracting away $\texttt{fact}$ from the term:
\[
  \texttt{fact}
  = \overbrace{\bigl(\lambda f.\,\lambda n.\; \texttt{if\ then\ else}\;(\texttt{iszero}\; n)\;(1)\;(\texttt{mult}\; n\;(f\;(\texttt{pred}\; n)))\bigr)}^{F}\;\texttt{fact}.
\]
\noindent Then the last equation becomes $\texttt{fact} = F \, \texttt{fact}$, which is a fixed point equation.
We can solve it up to $\beta$-equivalence by letting \( \texttt{fact} = \Theta F \). Using our Haskell program we define \texttt{fact}:
\begin{minted}{haskell}
  fact = a turing
      (l "f" (l "n"
          (a (a (a if_ (a iszero (v "n"))) (nat 1))
          (a (a mul (v "n"))
              (a (v "f")
              (a pred_ (v "n")))))
      ))
\end{minted}
\noindent We run the following command to compute $\texttt{fact} \,  \overline{\texttt{3}}$:
\begin{minted}{shell-session}
  ghci> normalize (a fact (nat 3))
\end{minted}
\noindent That of course yields $ \lambda f.\,\lambda x.\, f(f(f(f(f(f\,x))))) \;\equiv\; \overline{\texttt{6}}$ as result.
\begin{note}
  The latter implementation exhibits a striking syntactic resemblance to that of the \textbf{LISP} family of languages.
\end{note}

\subsubsection*{\centering Conclusions}
We have covered some of the most important notions that surround the untyped \lcalc. It has been shown to satisfy the Church-Rosser property of confluence, it has also been shown that this property is not the strongest of its kind. We have also seen that it is not strongly normalizing. For this reasons, it is in our best interest to add some spice to the \lcalc \ by means typing. This is to try and make it comply with some stronger requirements that will make it more practical to perform computation. We covered untyped \lcalc \ the next chapters focuses on minimal typing variations.
\begin{itemize}
\item \textbf{Untyped}: Where the type of expressions is never specified. As we have seen, this gives us maximal expressiveness at the cost of being unsafe. Informally, this means we are at a risk of composing unrelated terms that produce miscalculations. For instance $x M$.
\item \textbf{Minimal Typing}: Every expression is annotated with a type. This would be equivalent to the usual set theoretic approach to computation where functions have well defined domain and codomain.
\item \textbf{Extended Typing}: Similar to minimal typing but extended with constructions like products and sums.
\end{itemize}
This alternatives have dramatically different properties from one another. We will cover them in some detail throughout the text.


\paragraph{References and Resources.}{
  This chapter is influenced by many of the classic resources on the subject, primarily \cite{nederpelt2014:type-theory} for the basic aspects of the \lcalc. The proof to the Church-Rosser Theorem and the factorial example  are fine-tuned versions of what can be found in \cite{selinger2008:lambdanotes}, though, for what it's worth, the factorial implementation is original. Similar proofs are found in \cite{openlogicparallelbeta2025} and \cite{nagele2016mechanized}. Other text used as general resources are \cite{barendregt2000:intro} and \cite{girard1989:proofs-and-types}. We intentionally skipped some of the topics in the literature like combinatory algebras that are not that relevant to our focus of study.

  
}





